%==============================================================================
% ANEXA B - Cod JavaScript (Frontend)
% Contine toate fragmentele de cod frontend referite in Capitolele 2 si 3.
%==============================================================================

\chapter{Cod JavaScript (Frontend)}
\label{anexa:js}

Această anexă reunește fragmentele de cod JavaScript care implementează
componentele frontend descrise în Capitolele~2 și~3. Pentru concizie, unele
fragmente sunt prezentate parțial, păstrând doar partea reprezentativă;
codul integral este disponibil în repository-ul proiectului.

\section*{Configurarea hărții exterioare}

\begin{lstlisting}[language=JavaScript, caption={Parametrii de configurare ai hartii campusului}, label={lst:b1}]
const DEFAULT_CAMPUS_LAT = 44.433215;
const DEFAULT_CAMPUS_LON = 26.056764;
const DEFAULT_ZOOM_LEVEL = 18;
const MAP_BOUNDS = [
    [44.432102, 26.054500], // coltul de sud-vest
    [44.434329, 26.059028]  // coltul de nord-est
];
\end{lstlisting}

\section*{Construirea grafului și algoritmul Dijkstra (exterior)}

\begin{lstlisting}[language=JavaScript, caption={Construirea grafului din date GeoJSON}, label={lst:b2}]
function buildGraphFromGeoJSON(geoJSON) {
    const graph = {};
    const nodesMap = {};

    function addNode(lat, lng) {
        const key = getCoordKey(lat, lng);
        if (!graph[key]) {
            graph[key] = [];
            nodesMap[key] = L.latLng(lat, lng);
        }
        return key;
    }
    function addEdge(key1, key2, dist) {
        graph[key1].push({ node: key2, weight: dist });
        graph[key2].push({ node: key1, weight: dist });
    }
    // parcurgerea elementelor de tip linie din GeoJSON
    // si adaugarea nodurilor si muchiilor corespunzatoare
    return { graph, nodesMap };
}
\end{lstlisting}

\begin{lstlisting}[language=JavaScript, caption={Nucleul algoritmului lui Dijkstra}, label={lst:b3}]
function runDijkstra(graph, nodesMap, startKey, endKey) {
    const distances = {};
    const previous = {};
    const unvisited = new Set();

    for (const node in graph) {
        distances[node] = Infinity;
        previous[node] = null;
        unvisited.add(node);
    }
    distances[startKey] = 0;
    // se alege repetat nodul nevizitat cu distanta minima
    // si se actualizeaza distantele vecinilor
    // ...
    // reconstructia traseului din lantul de noduri precedente
}
\end{lstlisting}

\begin{lstlisting}[language=JavaScript, caption={Gasirea celui mai apropiat nod din graf}, label={lst:b4}]
function findNearestNode(targetLatLng, nodesMap) {
    let nearestKey = null;
    let minDistance = Infinity;
    for (const key in nodesMap) {
        const dist = targetLatLng.distanceTo(nodesMap[key]);
        if (dist < minDistance) {
            minDistance = dist;
            nearestKey = key;
        }
    }
    return nearestKey;
}
\end{lstlisting}

\section*{Securitatea la nivelul interfeței}

\begin{lstlisting}[language=JavaScript, caption={Prelucrarea textului inainte de afisare (anti XSS)}, label={lst:b5}]
function escapeHtml(text) {
    if (text == null) return '';
    const div = document.createElement('div');
    div.textContent = String(text);
    return div.innerHTML;
}
\end{lstlisting}

\section*{Indoor: persistare în baza de date}

\begin{lstlisting}[language=JavaScript, caption={Salvarea planului indoor curent in baza de date}, label={lst:b6}]
function saveIndoorDataToDatabase() {
    fetch('/api/indoor/save', {
        method: 'POST',
        headers: { 'Content-Type': 'application/json' },
        body: JSON.stringify(window.floorData)
    })
    .then(response => {
        if (!response.ok) throw new Error("Failed to save to database");
        return response.text();
    })
    .then(msg => {
        showCustomAlert("Plan salvat cu succes!");
    })
    .catch(error => {
        console.error("Error saving indoor data:", error);
        showCustomAlert("Eroare la salvarea planului.");
    });
}
\end{lstlisting}

\begin{lstlisting}[language=JavaScript, caption={Incarcarea planurilor indoor de pe server la pornire}, label={lst:b7}]
function loadIndoorData() {
    fetch('/api/indoor/load')
        .then(response => {
            if (!response.ok) throw new Error("Network response was not ok");
            return response.json();
        })
        .then(data => {
            if (data && Object.keys(data).length > 0) {
                window.floorData = data;
            }
        })
        .catch(error => {
            console.log("No existing indoor data found:", error);
            if (!window.floorData) window.floorData = {};
        });
}
\end{lstlisting}

\section*{Indoor: rutarea pe planul de etaj}

\begin{lstlisting}[language=JavaScript, caption={Conectarea nodurilor de start si destinatie la reteaua de coridoare}, label={lst:b8}]
function connectNodeToNetwork(nKey) {
    let p = nodesMap[nKey];
    let minDist = Infinity;
    let bestProj = null;
    let bestEdge = null;

    edgesList.forEach(edge => {
        if (nKey === edge.k1 || nKey === edge.k2) return;
        const res = distToSegment(p, edge.p1, edge.p2);
        if (res.dist < minDist) {
            minDist = res.dist;
            bestProj = res.proj;
            bestEdge = edge;
        }
    });

    // toleranta pentru liniile desenate "la ochi" care nu se ating exact in vertex
    if (bestEdge && (minDist < 25 || nKey === startKey || nKey === destKey)) {
        const projKey = addIndoorNode(bestProj.lat, bestProj.lng);
        // taiem muchia si inseram proiectia, apoi conectam punctul nostru
        graph[bestEdge.k1].push({ node: projKey, weight: euclidean(bestEdge.p1, bestProj) });
        graph[projKey].push({ node: bestEdge.k1, weight: euclidean(bestEdge.p1, bestProj) });
        graph[bestEdge.k2].push({ node: projKey, weight: euclidean(bestEdge.p2, bestProj) });
        graph[projKey].push({ node: bestEdge.k2, weight: euclidean(bestEdge.p2, bestProj) });
        graph[nKey].push({ node: projKey, weight: euclidean(p, bestProj) });
        graph[projKey].push({ node: nKey, weight: euclidean(p, bestProj) });
    }
}
\end{lstlisting}
